% Options for packages loaded elsewhere
\PassOptionsToPackage{unicode}{hyperref}
\PassOptionsToPackage{hyphens}{url}
\PassOptionsToPackage{dvipsnames,svgnames,x11names}{xcolor}
%
\documentclass[
  11pt,
]{article}
\usepackage{amsmath,amssymb}
\usepackage{iftex}
\ifPDFTeX
  \usepackage[T1]{fontenc}
  \usepackage[utf8]{inputenc}
  \usepackage{textcomp} % provide euro and other symbols
\else % if luatex or xetex
  \usepackage{unicode-math} % this also loads fontspec
  \defaultfontfeatures{Scale=MatchLowercase}
  \defaultfontfeatures[\rmfamily]{Ligatures=TeX,Scale=1}
\fi
\usepackage{lmodern}
\ifPDFTeX\else
  % xetex/luatex font selection
\fi
% Use upquote if available, for straight quotes in verbatim environments
\IfFileExists{upquote.sty}{\usepackage{upquote}}{}
\IfFileExists{microtype.sty}{% use microtype if available
  \usepackage[]{microtype}
  \UseMicrotypeSet[protrusion]{basicmath} % disable protrusion for tt fonts
}{}
\makeatletter
\@ifundefined{KOMAClassName}{% if non-KOMA class
  \IfFileExists{parskip.sty}{%
    \usepackage{parskip}
  }{% else
    \setlength{\parindent}{0pt}
    \setlength{\parskip}{6pt plus 2pt minus 1pt}}
}{% if KOMA class
  \KOMAoptions{parskip=half}}
\makeatother
\usepackage{xcolor}
\usepackage[margin=1in]{geometry}
\setlength{\emergencystretch}{3em} % prevent overfull lines
\providecommand{\tightlist}{%
  \setlength{\itemsep}{0pt}\setlength{\parskip}{0pt}}
\setcounter{secnumdepth}{5}
\usepackage{amsmath,amssymb,amsthm,mathtools}
\newcommand{\ri}{\operatorname{ri}}
\newcommand{\rec}{\operatorname{rec}}
\newcommand{\conv}{\operatorname{conv}}
\newcommand{\cone}{\operatorname{cone}}
\newcommand{\epi}{\operatorname{epi}}
\newcommand{\dist}{\operatorname{dist}}
\newcommand{\Span}{\operatorname{span}}
\newcommand{\covol}{\operatorname{covol}}
\ifLuaTeX
  \usepackage{selnolig}  % disable illegal ligatures
\fi
\usepackage{bookmark}
\IfFileExists{xurl.sty}{\usepackage{xurl}}{} % add URL line breaks if available
\urlstyle{same}
\hypersetup{
  pdftitle={Finite balanced complexes: uniform cellwise slopes and exact single-cell real polyhedral majorants},
  pdfauthor={Independent scoped verification draft},
  colorlinks=true,
  linkcolor={Maroon},
  filecolor={Maroon},
  citecolor={Blue},
  urlcolor={Blue},
  pdfcreator={LaTeX via pandoc}}

\title{Finite balanced complexes: uniform cellwise slopes and exact
single-cell real polyhedral majorants}
\author{Independent scoped verification draft}
\date{8 October 2026}

\begin{document}
\maketitle
\begin{abstract}
We independently reconstruct the six units of working299, the five units
of working296, and the three units of working297. On a finite globally
embedded rational complex with strictly positive actual real balanced
top weights and no omitted boundary, a continuous
printed-polyhedrally-concave function within a bounded distance of a
cellwise affine real function has uniform Euclidean Lipschitz bounds on
every closed cell. At every face point it admits a uniformly bounded
ambient affine support on the entire incident cells. The proof uses
rational feasible-face approximation, an explicit polyhedral distance
estimate, covolume conversion of balancing into stationary flux, an
extended-real directional-derivative argument, and induction over all
faces. On one closed rational cell, the resulting Lipschitz hypothesis
permits an exact convex extension to the whole affine hull and
rational-carrier real piecewise-affine majorants with actual recession
and prescribed finite rational contacts. An unbalanced orthant example
has optimal uniform approximation error one half for every finite real
piecewise-affine competitor. These are scoped results. Shared-face trace
matching and common-support compatibility after approximation remain
unproved; the original P05.1 and P05.2 classifications remain OPEN. This
document is an unverified draft, with candidate=false.
\end{abstract}

{
\hypersetup{linkcolor=}
\setcounter{tocdepth}{3}
\tableofcontents
}
\section{Status, domain, and the exact
assertions}\label{status-domain-and-the-exact-assertions}

The inputs working296, working297, and working299 are proposals rather
than grants. We verify them below without importing their proofs as
accepted results. The six supplied whole grants are A224, B225, R230,
S262, S272, and S291, with their supplied source and reader
qualifications. They are used only at their stated scope. In particular,
the main argument does not invoke a compact maximum principle on a
noncompact space. It does not invoke an operator on a lone unbalanced
cell, or on a non-piecewise-affine intermediate extension. The detailed
source ledger is a separate artifact. Source quotations and source
pixels have not been incorporated in this document.

Let \(X\) be a finite pure \(d\)-dimensional rational polyhedral complex
globally embedded in a finite-dimensional Euclidean space with a fixed
lattice. All its closed faces are retained. Distinct cells meet in
common faces. Fix a rational refinement of this complex. For every top
cell \(\sigma\), its lattice is the saturated lattice in
\(T_\sigma=\operatorname{span}(\sigma-\sigma)\). The top weights
\(w_\sigma\) are strictly positive actual real numbers. At every
codimension-one face \(\rho\), balancing means \[
\sum_{\sigma\succ\rho} w_\sigma u_{\sigma/\rho}=0
\quad\hbox{in the ambient quotient by }T_\rho,
\] where \(u_{\sigma/\rho}\) is the primitive quotient lattice vector
directed into \(\sigma\). Balancing is required at every such face,
including every face which would otherwise be a boundary. No rationality
of the weights is assumed.

A \emph{printed PCC} is the following finite inequality. At a rational
center \(y\) in a face \(\tau\), choose rational endpoints \(x_i\) in
assigned cells \(\sigma_i\succeq\tau\) and nonnegative real coefficients
\(\alpha_i\) such that \(\sum_i\alpha_i=1\) and
\(\sum_i\alpha_i x_i=y\). The inequality is
\[ f(y)\leq\sum_i\alpha_i f(x_i). \] The printed-PC hypothesis means
that these inequalities hold for the rational structures and incidences
in the definition. We will use the fixed refinement. Continuity of \(f\)
is a separate, indispensable hypothesis. We do not assume that an
arbitrary function satisfying only the rational-point inequalities is
continuous.

Let \(\gamma\) be an actual continuous real piecewise-affine function
affine on each cell of the fixed refinement. Put \[
G=\max_\sigma\|a_\sigma\|,
\qquad \gamma|_\sigma(x)=\langle a_\sigma,x\rangle+b_\sigma,
\quad a_\sigma\in T_\sigma.
\] Here \(a_\sigma\) is the intrinsic gradient, orthogonally embedded in
the ambient space. Faces inherit the same bound. Neither convexity nor
printed-PC nor positive growth of \(\gamma\) is needed in the main
assertion. We take \(C\geq0\). If the complex is empty, all conclusions
are vacuous and we set \(G=A_X=B_X=0\); henceforth the main proof treats
a nonempty complex.

\textbf{Theorem 1 (working299, all six units).} There are finite
constants \(A_X,B_X\) depending on the fixed cell geometry and its
Euclidean metric, and not on \(f\), its location, \(C\), or the values
of \(\gamma\), with the following property. If \(f\) is continuous and
printed-PC and \(|f-\gamma|\leq C\) on \(X\), then every closed convex
cell has global Euclidean Lipschitz constant at most \(A_XC+B_XG\). At
every \(y\in X\) there is an ambient vector \(p_y\) of norm at most the
same bound such that \[
f(x)\geq f(y)+\langle p_y,x-y\rangle
\quad\hbox{for every }x\hbox{ in every cell incident to the minimal face of }y.
\] In particular, the support extends over entire incident cells. The
constants can be enlarged once so that both conclusions use the same
pair \(A_X,B_X\).

The metric conclusion is cellwise, and consequently bounds variation
along any finite polygonal path whose segments lie in cells by its
Euclidean length. It is not an unqualified ambient-distance statement
across nonincident cells or components. For example, two disjoint
parallel affine lines, each carrying a positive balanced one-dimensional
weight, can have different affine slopes. The ambient distance between
corresponding points stays fixed while their function values separate
linearly. Their union meets the main hypotheses with \(C=0\).

\textbf{Theorem 2 (working296, all five units).} Let \(\sigma\) be a
nonempty closed rational convex cell in its rational affine hull \(Q\).
Let \(\ell\) be an actual real affine function on \(Q\) whose slope is
strictly positive on every nonzero direction in
\(\operatorname{rec}(\sigma)\). Suppose \(f\) is finite continuous
convex on \(\sigma\), \(|f-\ell|\leq C\), and \(f\) has a finite global
Euclidean Lipschitz constant. There is a proper convex real
piecewise-affine auxiliary function \(\Gamma\) on all of \(Q\), with
rational spatial carrier, equal to \(\ell\) on all of \(\sigma\). There
is a finite continuous convex function \(F\) on all of \(Q\), equal to
\(f\) on \(\sigma\), with \(|F-\Gamma|\leq C\). For every \(\eta>0\) and
every prescribed finite set of rational points of \(\sigma\), there is a
genuine finite convex real piecewise-affine majorant \(P\) on \(Q\),
with rational spatial carrier, such that \(F\leq P\leq F+\eta\), with
equality at the prescribed points and with \(P^\infty=\Gamma^\infty\).
Restriction gives uniform above-approximants on \(\sigma\) with the
actual cell recession. A common dyadic shift gives uniformly decreasing
approximants. Exact contacts refer to the unshifted construction.

When \(f\) is as in Theorem 1 and \(\ell=\gamma|_\sigma\) for an
admitted positive-growth real PA PC reference, the Lipschitz assumption
in Theorem 2 is supplied by Theorem 1. The auxiliary \(\Gamma\) does not
replace the actual reference on \(X\).

\textbf{Theorem 3 (working297, all three units).} On the unbalanced
orthant \(\sigma=\mathbb R_{\geq0}^2\), let \[
\gamma(x,y)=x+y,
\qquad f(x,y)=x+y+\max(0,1-\sqrt{xy}).
\] This is continuous convex, differs from \(\gamma\) by a number in
\([0,1]\), and has the same recession. Every finite continuous real PA
function \(a\) on this orthant satisfies \(\|a-f\|_\infty\geq\tfrac12\).
The bound is attained by \(a=\gamma+\tfrac12\). This is a boundary
diagnostic outside the balanced domain, not a counterexample to Theorem
1 or to the original balanced task.

\section{Elementary foundations used in the
reconstruction}\label{elementary-foundations-used-in-the-reconstruction}

We disclose the following ordinary foundations: completeness of the real
numbers; compactness of closed bounded subsets of finite-dimensional
Euclidean space; elementary Gaussian elimination and finite-dimensional
linear algebra; the product rule and divergence theorem on polyhedral
domains with spherical truncation; and one-variable convex secant
inequalities. Here are the additional consequences we need, with proofs
to specify their scope.

\textbf{Finite cones and a coefficient bound.} For a fixed finite list
\(v_1,\ldots,v_m\), every vector in their nonnegative cone has a
representation using a linearly independent sublist. Indeed, if positive
coefficients have dependent columns, take a nonzero dependence
\(\sum_i b_i v_i=0\), reversing its sign to have a positive entry.
Subtract \(s b_i\) from the coefficients, where
\(s=\min_{b_i>0} a_i/b_i\). Coefficients remain nonnegative, a positive
coefficient becomes zero, and the represented vector does not change.
Repeat. For each independent sublist \(I\), its column map has a linear
inverse on its image. Thus \[ \sum_{i\in I} a_i\leq K_I\|v\|. \] There
are only finitely many sublists; their maximum \(K\) works for every
vector in the cone. This also proves that the cone is closed: for a
convergent sequence of represented vectors choose a subsequence with one
independent sublist; its coefficients are bounded and converge along a
further subsequence. The empty cone is \(\{0\}\) and has coefficient
bound zero.

If a point \(w\) is outside a closed convex cone \(K\), nearest-point
separation is elementary. A nearest point \(k\in K\) exists by
restricting the distance minimization to a sufficiently large closed
ball. First variation gives \(\langle w-k,z-k\rangle\leq0\) for
\(z\in K\). Scaling \(k\) and using \(0\in K\) gives
\(\langle w-k,k\rangle=0\). Consequently \(v=w-k\) satisfies
\(\langle v,z\rangle\leq0\) on \(K\) and
\(\langle v,w\rangle=\|v\|^2>0\). This is the only cone separation
needed below.

\textbf{Finite polyhedral generators and projections.} A rational cone
given by finitely many homogeneous inequalities has a rational basis for
its lineality space, obtained from the joint kernel of its rows. Split
off this space using a rational complementary subspace. In the pointed
quotient, orient the rows as nonnegative functions on the cone. Their
sum is positive on every nonzero vector; otherwise all rows vanish,
giving lineality. Its section at value one is bounded: an unbounded
sequence, normalized and passed to a unit limit, would give a nonzero
cone vector at value zero. A bounded polyhedron is the convex hull of
finitely many vertices. To see this without importing a representation
theorem, at a nonvertex point choose a nonzero direction annihilating
its active rows. Moving in both signs reaches additional active rows,
since the polyhedron is bounded. The original point is a convex
combination of the two endpoints. Induction on active-row rank ends at
vertices. Vertices solve finitely many independent systems. Rational
systems have rational solutions. This proves finite rational generation
of the pointed cone; adjoin both signs of a rational lineality basis.

Homogenizing a nonempty polyhedron \(A x\leq b\) as \(A t\leq\alpha b\),
\(\alpha\geq0\), gives, by the same cone result,
\[ P=\operatorname{conv}\{s_1,\ldots,s_r\}+\operatorname{cone}\{r_1,\ldots,r_k\}. \]
Generators with positive \(\alpha\) yield the \(s_i\); those with zero
\(\alpha\) yield recession directions. All spatial data are rational
when the original rows are rational. Conversely, a finite convex hull
plus a finite cone is a polyhedron: express it using nonnegative
coefficients, a sum-one equality for hull coefficients, and a linear
image. Eliminate variables one at a time. For one variable, pair every
lower bound with every upper bound and retain rows not involving that
variable. These finite weak inequalities are exactly the condition that
the interval of allowable values is nonempty; if one side is absent the
interval is unbounded and still nonempty. Equalities are two
inequalities. This proves finite elimination, over either real or
rational coefficients, and rationality of the projected polyhedron when
the spatial data are rational. These arguments also apply in a relative
affine hull.

\textbf{Local convex bounds and relative subgradients.} A finite convex
function \(u\) on a relatively open neighborhood of \(q\) has local
Lipschitz bounds, hence continuity, and a relative subgradient at \(q\).
Choose a closed cube centered at \(q\) inside that neighborhood.
Convexity bounds \(u\) above by the largest vertex value. Reflection
about \(q\) gives a lower bound on a smaller cube. On a still smaller
cube, extend the segment between any two points a fixed positive
distance in either direction within the bounded cube; the convex secant
inequalities bound its slope by the oscillation divided by that
distance. This proves local Lipschitzness without presupposing
continuity.

For existence of a subgradient, use the closed convex epigraph over this
cube. Project \((q,u(q)-t)\) to it, and normalize the vector from the
projection to that point. As \(t\downarrow0\) the projections converge
to \((q,u(q))\); compactness of the unit sphere supplies a limiting
nonzero supporting normal. Its vertical component is nonpositive because
the epigraph contains upward rays. It cannot be zero: a nonzero purely
spatial functional cannot have its maximum over a spatial neighborhood
at its center. Dividing by the negative vertical component gives an
affine lower support. Convex secants extend this local support over the
entire original convex domain. A zero-dimensional relative domain has
the zero subgradient. The same argument on full affine space proves that
a finite continuous convex function is the supremum of all its affine
supports. In particular its double convex conjugate equals itself. No
general duality theorem is being imported here.

\section{L0: real endpoint PCC and convexity on inherited closed
cells}\label{l0-real-endpoint-pcc-and-convexity-on-inherited-closed-cells}

Fix incident cells \(\sigma_i\succeq\tau\) and a real PCC configuration
with \(\alpha_i>0\); delete zero coefficients first. Write perspective
variables \[ (\alpha_i,t_i),\qquad t_i=\alpha_i x_i. \] The conditions
\(x_i\in\sigma_i\), \(\sum\alpha_i=1\), and \(\sum t_i\in\tau\) are
finitely many rational weak linear inequalities and equalities in these
variables. For example, \(A_i x_i\leq b_i\) becomes
\(A_i t_i\leq\alpha_i b_i\). Equations for relative affine hulls are
included. Add \(\alpha_i\geq0\).

At the given configuration, impose as equations every inequality that is
active. This is its minimal feasible face. Gaussian elimination gives a
rational particular solution and a rational basis for the homogeneous
solution space. Rational combinations are dense in that affine space.
All inactive rows have strictly positive slack at the given point, and
there are only finitely many; sufficiently close approximations preserve
these inequalities. They also preserve \(\alpha_i>0\). Thus the given
configuration is approximated by rational feasible configurations. Their
endpoints \(t_i/\alpha_i\) and centers \(\sum_i t_i\) are rational, and
their coefficients are permitted printed PCC coefficients. A center in
the relative interior stays there under sufficiently close feasible
approximations; for a boundary center use its minimal inherited face,
which is contained in every assigned incident cell. Apply printed PC
there and take limits using continuity of \(f\). This proves PCC for
real endpoints, real centers, and real coefficients in every fixed face
incidence.

For two real points in one inherited closed cell \(\sigma\), their
convex combination belongs to a unique minimal face
\(\tau\preceq\sigma\). Assign both endpoints to \(\sigma\), and apply
the just-proved real PCC at that center. This proves convexity of \(f\)
on \(\sigma\). No face is merged or coarsened. In particular, convexity
on an inherited closed face follows without any automatic-continuity
assertion about raw rational-point PC functions.

\section{L1: quantitative distance and uniform
insets}\label{l1-quantitative-distance-and-uniform-insets}

Let \(P=\{x:n_j\cdot x\leq b_j\}\) be nonempty in a fixed Euclidean
affine hull, with equality constraints already removed by working in
that hull. A nearest point \(z\) to \(x\) exists by compactness in a
sufficiently large ball. Put \(w=x-z\), and let \(J\) be the active rows
at \(z\). We claim \[ w\in\operatorname{cone}\{n_j:j\in J\}. \]
Otherwise cone separation proved above gives \(v\) with
\(n_j\cdot v\leq0\) for all active rows and \(w\cdot v>0\). For
sufficiently small \(s>0\), \(z+s v\) satisfies the inactive rows, by
their positive slack, and also every active row. Nearest-point first
variation would require \(w\cdot v\leq0\), a contradiction.

Independent-column elimination gives \(w=\sum_{j\in J}\lambda_j n_j\)
with \(\lambda_j\geq0\) and \(\sum\lambda_j\leq H\|w\|\), where \(H\)
depends only on the finite row list, not on the right-hand sides or
active subset. Hence, for \(\delta(x)=\max(0,\max_j(n_j\cdot x-b_j))\),
\[
\|w\|^2=\sum_j\lambda_j(n_j\cdot x-b_j)
\leq H\delta(x)\|w\|,
\qquad \operatorname{dist}(x,P)\leq H\delta(x).
\] Zero rows cause no difficulty; a nonempty system has their right-hand
sides nonnegative, and they never help a nonzero representation. An
empty row list gives the whole affine hull and distance zero. In
dimension zero the assertion is also immediate. This is a proof of the
needed distance estimate, including its uniformity; the name of a
distance theorem is not a substitute for the argument.

For a face \(\tau\), list all its own facet inequalities, together with
all inequalities of every incident top cell whose restrictions have
strictly positive slack on \(\operatorname{ri}(\tau)\). Constraints
identically equal on \(\tau\) are not strengthened. A valid affine
inequality on \(\tau\) which vanishes at one relative interior point
vanishes everywhere on \(\tau\); otherwise its affine restriction would
change sign in a small relative neighborhood. This explains the
dichotomy and the finite list.

Choose a rational point \(z_0\in\operatorname{ri}(\tau)\), and a
rational margin \(\delta_\tau>0\) below every listed slack at \(z_0\).
Strengthen every listed row by that margin, retaining the other defining
constraints of \(\tau\). The resulting inset \(\tau^\delta\) is
nonempty. For every \(y\in\tau\), the violation of each strengthened row
is at most \(\delta_\tau\). The preceding distance estimate gives a
nearest point \(z\in\tau^\delta\) with
\[ \|z-y\|\leq D_\tau:=H_\tau\delta_\tau. \] This estimate holds
uniformly near lower faces and at infinity. A fixed interior anchor
alone would give no such uniform bound. If the list is empty, take
\(\tau^\delta=\tau\) and \(D_\tau=0\). A zero-dimensional face also has
\(z=y\) and \(D_\tau=0\), even when some listed restrictions are
positive constants.

For any finite list of assigned tangent-cone generators \(v_j\), choose
one rational \(\varepsilon_\tau>0\) sufficiently small that
\[ z+\varepsilon_\tau v_j\in\sigma(j)
\quad\hbox{for every }z\in\tau^\delta.
\] Indeed, an incident-cell row active on \(\tau\) is respected by the
assigned cone direction. Every other row has slack at least
\(\delta_\tau\) on the inset, and its possible positive change is
bounded by \(\varepsilon_\tau\max_j|n\cdot v_j|\). Affine-hull
equalities are preserved. Finiteness supplies one choice for all rows
and generators. Empty-list and zero-dimensional cases are included by
the same reasoning.

\section{L2: positive real balancing fills the local
span}\label{l2-positive-real-balancing-fills-the-local-span}

At \(y\in\operatorname{ri}(\tau)\) the tangent cones \(C_{\sigma,\tau}\)
of incident top cells are independent of the choice of \(y\): they are
given by the rows active identically on \(\tau\), with lineality
containing \(T_\tau\). They are rational polyhedral cones of dimension
\(d\). By the elementary generation result, choose finite rational
generators for them, including both signs of a rational basis of
\(T_\tau\), each assigned to an incident top cell. Put
\[ E_\tau=\operatorname{span}\bigcup_{\sigma\succeq\tau} C_{\sigma,\tau}. \]

The star inherits balancing. A facet of a tangent cone comes from a
facet of its cell containing \(\tau\), the adjacent star cells are
exactly the adjacent original cells there, and their tangent lattices
and quotient primitive normals are unchanged. A rational refinement
preserves balancing as well: a new facet in the interior of an old cell
has two opposite primitive quotient directions with equal inherited
weight; an old facet is subdivided with its old relation unchanged. This
specifies rather than assumes the local inheritance used here.

For a facet \(\rho\prec\sigma\), write
\(c_\sigma=\operatorname{covol}(N_\sigma)\) and
\(c_\rho=\operatorname{covol}(N_\rho)\) in their Euclidean spans.
Saturation implies \(N_\sigma/N_\rho\) is free of rank one. A lattice
basis for \(N_\rho\), together with a lift \(u\) of its primitive
quotient generator, is a basis for \(N_\sigma\). The parallelepiped
volume identity therefore gives \[ \|u_\perp\|=c_\sigma/c_\rho,
\qquad u_\perp=(c_\sigma/c_\rho)\nu_{\sigma/\rho}, \] where
\(\nu_{\sigma/\rho}\) is the inward Euclidean unit conormal.
Consequently actual real balancing is exactly the stationary
unit-conormal flux relation \[
\sum_{\sigma\succ\rho}\rho_\sigma\nu_{\sigma/\rho}=0,
\qquad \rho_\sigma=w_\sigma c_\sigma>0.
\] This covolume conversion is essential: replacing primitive normals by
unit normals without it is generally incorrect.

Let \(K\) be the conical hull of the finite generators in \(E_\tau\). It
is closed. Suppose \(K\ne E_\tau\). The cone separation argument
supplies a nonzero linear functional \(\lambda\) on \(E_\tau\) that is
nonnegative on \(K\). It is positive somewhere on the union of the
cones, since that union spans \(E_\tau\).

Apply the divergence theorem on each \(C_{\sigma,\tau}\cap B(0,1)\) to
the constant vector obtained by projecting \(\nabla\lambda\) to
\(T_\sigma\), with density \(\rho_\sigma\). Its divergence is zero. The
sum of planar facet fluxes cancels by the preceding stationary relation,
with outward conormals the negatives of the displayed inward conormals.
The remaining spherical flux is \[
\sum_\sigma \rho_\sigma
\int_{C_{\sigma,\tau}\cap S^{d-1}_{T_\sigma}}\lambda(\theta)\,dS(\theta).
\] Every term is nonnegative. One is strictly positive: if \(\lambda\)
is positive on a point of a full-dimensional cone, it is positive on a
relatively open spherical patch of that cone. Lower-dimensional edges
have zero boundary measure. This contradicts zero total flux.

In dimension one the same argument is the two-endpoint divergence
formula on intervals or rays, with counting measure on the unit sphere.
A whole line has canceling antipodal endpoints. Cones with lineality are
handled by the same ball truncation; no pointedness assumption is
needed. If \(T_\tau=E_\tau\), its two-sided tangent basis already proves
the conclusion. In dimension zero, \(E_\tau=\{0\}\). These observations
cover zero relative codimension and the degenerate cases without a
complete-plane hypothesis.

We have proved \(K=E_\tau\). Independent-column elimination now gives
one finite geometric constant \(K_\tau\) such that both \(v\) and \(-v\)
admit assigned nonnegative generator representations with coefficient
sums at most \(K_\tau\|v\|\). The constant is calculated from finitely
many inverse column maps. Actual weight magnitudes do not enter that
bound; their positivity and balancing were needed to prove that the cone
fills the span.

\section{L3: finite directional costs and supports on whole incident
cells}\label{l3-finite-directional-costs-and-supports-on-whole-incident-cells}

Assume \(f|_\tau\) has global Lipschitz constant \(M_\tau\). For a point
face set \(M_\tau=0\). Fix \(y\in\operatorname{ri}(\tau)\) and choose
the inset point \(z\) of L1. Write \(\xi=z-y\), so \(\xi\in T_\tau\) and
\(\|\xi\|\leq D_\tau\). For a feasible direction
\(w\in C_{\sigma,\tau}\), define
\[ h_\sigma(w)=\lim_{t\downarrow0}\frac{f(y+t w)-f(y)}t. \] Convex
secants show that this limit exists in \([-\infty,\infty)\): a fixed
feasible positive secant bounds it above, but initially it may equal
\(-\infty\). Positive homogeneity follows by a change of variable.
Midpoint convexity, applied to \(y+2t u\) and \(y+2t v\), shows
\[ h_\sigma(u+v)\leq h_\sigma(u)+h_\sigma(v). \] This extended-real
inequality is legitimate before finiteness: no term can be \(+\infty\),
and if one right-hand term tends to \(-\infty\), the midpoint bound
forces the corresponding upper limit to be \(-\infty\). Tangent
directions have common derivatives in all incident cells, since a small
segment in either tangent direction stays in \(\tau\). In particular
\(h_\sigma(-\xi)\leq M_\tau D_\tau\) and is finite.

For an assigned generator \(v_j\), L1 puts
\(y+\xi+\varepsilon v_j=z+\varepsilon v_j\) in its cell \(\sigma(j)\).
The bound to the \emph{same} affine restriction of \(\gamma\) gives
\[ h_{\sigma(j)}(\xi+\varepsilon v_j)
\leq 2C+G(D_\tau+\varepsilon\|v_j\|). \] Subadditivity applied to
\(\varepsilon v_j=(\xi+\varepsilon v_j)+(-\xi)\) gives the uniform
generator upper bound \[
h_{\sigma(j)}(v_j)\leq U_\tau,
\qquad U_\tau=\frac{2C+(G+M_\tau)D_\tau}{\varepsilon_\tau}+G V_\tau,
\quad V_\tau=\max_j\|v_j\|.
\] An upper bound here allows a value \(-\infty\); we remove that
possibility next, instead of assuming it away.

Any finite list of assigned feasible directions \(w_a\) with
\(\sum_a w_a=0\) has feasible endpoints \(y+t w_a\) for one common
sufficiently small \(t>0\). Their equal-weight barycenter is \(y\). The
real-endpoint PCC of L0 gives \[ \sum_a\frac{f(y+t w_a)-f(y)}t\geq0. \]
If an incident derivative \(h_\sigma(w)\) were \(-\infty\), append a
fixed-generator representation of \(-w\) from L2, using its nonnegative
coefficients as scalings of the assigned directions and deleting zeros.
Every appended directional quotient has a finite upper limit, bounded by
its coefficient times \(U_\tau\). The sum would tend to \(-\infty\),
contradicting the displayed PCC. Thus every incident directional
derivative is finite. Taking limits now shows that every finite zero-sum
assigned list has nonnegative total derivative cost.

Define, over \emph{all} finite assigned incident-direction
decompositions, \[
H_y(v)=\inf\left\{\sum_a h_{\sigma(a)}(w_a):
w_a\in C_{\sigma(a),\tau},\ \sum_a w_a=v\right\}.
\] The empty decomposition has cost zero at \(v=0\). The generator
representation gives \(H_y(v)\leq K_\tau U_\tau\|v\|\). Append a
generator representation of \(-v\) to any candidate decomposition of
\(v\); the zero-sum cost inequality gives the opposite bound.
Consequently
\[ -K_\tau U_\tau\|v\|\leq H_y(v)\leq K_\tau U_\tau\|v\|. \] There is no
\(-\infty\) infimum. Scaling decompositions proves positive homogeneity
for positive scalars; the zero-sum property proves \(H_y(0)=0\).
Concatenating two decompositions whose costs are within a positive
tolerance of their finite infima proves subadditivity. These properties
and the two-sided bound imply that \(H_y\) is globally Lipschitz with
constant \(K_\tau U_\tau\).

Apply the full-space epigraph-support construction from the foundations
to this finite continuous sublinear function at zero. It gives a linear
functional \(p_y\leq H_y\). Testing both \(v\) and \(-v\) gives
\(\|p_y\|\leq K_\tau U_\tau\). Extend it by zero on \(E_\tau^\perp\) to
an ambient vector. For every \(x\) in an incident cell,
\(x-y\in C_{\sigma,\tau}\) and convex secants give \[
f(x)-f(y)\geq h_\sigma(x-y)\geq H_y(x-y)\geq\langle p_y,x-y\rangle.
\] The support therefore holds on the entire closed incident cell,
including its far end and its boundary, rather than only on a germ. When
\(E_\tau=\{0\}\) take \(p_y=0\), omitting all generator maxima and
inverse constants. A zero-dimensional face in a positive-dimensional
complex still has a nonzero local span in general and uses the full
construction. All uses of real PCC, extended-real limits, subadditivity,
and supports have now been justified.

\section{L4 and L5: induction over every face and the quantitative
conclusion}\label{l4-and-l5-induction-over-every-face-and-the-quantitative-conclusion}

Order all the finitely many faces by dimension. Before handling a face
\(\tau\), supports have been constructed at all points of its proper
faces, by using each point's minimal face. Let \(B\) be the maximum of
their norm bounds, or zero if there are no proper faces. These supports
hold on all of \(\tau\): any top cell containing \(\tau\) is incident to
the minimal face of each such boundary point.

At \(q\in\operatorname{ri}(\tau)\), relative subgradients exist by the
foundations. Let \(a\) be \emph{any} such subgradient. If \(a\ne0\), put
\(e=a/\|a\|\) and consider \(q+t e\). Because \(q\) is relatively
interior, it is feasible for a positive initial interval. If the ray
exits at a finite \(t_0\), its endpoint \(z\) belongs to a proper face.
The subgradient at \(q\) and the already constructed support at \(z\)
imply \[
t_0\|a\|\leq f(z)-f(q)\leq t_0\langle p_z,e\rangle\leq t_0 B.
\] If it never exits, bounded difference from the affine restriction
gives \[
\lim_{t\to\infty}\frac{f(q+t e)-f(q)}t=\langle a_\gamma,e\rangle\leq G,
\] while the subgradient inequality bounds every such secant below by
\(\|a\|\). Thus every relative subgradient has norm at most
\(M_\tau=\max(B,G)\). The zero subgradient already satisfies this bound.
For a zero-dimensional face use the sharper \(M_\tau=0\).

Supports at two relative interior points give the Lipschitz inequality
there. To pass to arbitrary points of the closed face, mix each point
with a fixed relative interior point using the same positive mixing
parameter, then let that parameter decrease to zero and use continuity.
This proves \emph{global} Lipschitzness on the whole possibly unbounded
face. L3 now supplies its uniform ambient support bound
\[ B_\tau=K_\tau\left(\frac{2C+(G+M_\tau)D_\tau}{\varepsilon_\tau}+G V_\tau\right). \]

This is induction over all finite faces, not a vertex-started induction.
A minimal positive-dimensional face has no proper faces; every relevant
ray then remains forever, so the recession argument starts the
induction. Lineality presents no problem. An affine top cell with no
proper faces is also included. Positivity of \(f\), positive growth of
\(\gamma\), averaging of references, independence of quotient normals,
and a complete-plane assumption are absent from the argument.

For explicit dependence, suppose proper-face support bounds have
coefficients \(A_\rho C+B_\rho G\). Set
\[ a_\tau=\max_{\rho\prec\tau} A_\rho,\quad
b_\tau=1+\max_{\rho\prec\tau}B_\rho, \] with empty maxima zero. Then
\(M_\tau\leq a_\tau C+b_\tau G\). One admissible next support
coefficient pair is \[
A_\tau=K_\tau\left(2/\varepsilon_\tau+a_\tau D_\tau/\varepsilon_\tau\right),
\quad B_\tau^{\rm coeff}=K_\tau\left((1+b_\tau)D_\tau/\varepsilon_\tau+V_\tau\right).
\] Zero-dimensional faces have Lipschitz bound zero but may have nonzero
ambient support bounds obtained from L3 when the ambient local span is
nonzero. Thus one must not set their ambient support bound to zero
unless the entire local span is zero. This distinction is retained in
the recursion. Taking maxima of the finitely many Lipschitz and support
coefficient pairs proves Theorem 1. Summing the cellwise inequality on
the segments of an admissible polygonal path proves its stated intrinsic
consequence. This completes L0--L5 of working299; the regularization
input supplied by L5 is a cellwise input and not a gluing theorem.

\section{C1 and C2: a rational-wall auxiliary reference on one affine
hull}\label{c1-and-c2-a-rational-wall-auxiliary-reference-on-one-affine-hull}

Identify the rational affine hull \(Q\) with a Euclidean affine space
and write \[ \sigma=\{x\in Q:A_j\cdot x\leq b_j\},\qquad
\ell(x)=a\cdot x+b. \] All constraint rows are rational; the slope \(a\)
and constant \(b\) are actual reals. Put
\[ \delta(x)=\max(0,\max_j(A_j\cdot x-b_j)),\qquad
\Gamma(x)=\ell(x)+\kappa\delta(x), \] where \(\kappa\) is a positive
integer to be chosen. The violation function is convex PA with rational
walls. Adding the same real affine function to all its pieces leaves
those walls rational. Thus \(\Gamma\) is a genuine scalar real PA
function, exactly equal to \(\ell\) on the entire cell.

Let \(r(v)=\max(0,\max_j A_j\cdot v)\). The recession function of
\(\Gamma\) is \[ h(v)=a\cdot v+\kappa r(v). \] On the compact set
\(S=\operatorname{rec}(\sigma)\cap\{\|v\|=1\}\), the assumed slope has a
positive minimum if \(S\) is nonempty. It remains positive in a
neighborhood of \(S\) on the unit sphere. On the compact complement of a
smaller such neighborhood, \(r\) has a positive minimum; a sufficiently
large integer \(\kappa\) makes \(h\) positive there as well. If \(S\) is
empty, \(r\) is positive everywhere on the unit sphere and the same
conclusion follows directly. Hence \(h(v)\geq c\|v\|\) for some \(c>0\).
Since \(\Gamma-h\) is uniformly bounded by the finitely many constant
terms in its affine pieces, \(\Gamma(x)\geq c\|x\|-D\) after choosing an
origin. This proves proper coercive growth. The zero-dimensional case is
a constant on a point and is proper, with the recession condition
vacuous. A full positive-dimensional affine space cannot satisfy the
slope assumption on both \(v\) and \(-v\), so there is no overlooked
full-space example contradicting that assumption.

The L1 distance proof applies verbatim to this cell, with a finite
constant \(H\) and \(\operatorname{dist}(x,\sigma)\leq H\delta(x)\).
This independently verifies C1. The preceding construction verifies C2.
It is an auxiliary extension only; neither its values off \(\sigma\) nor
its coefficients are substituted for the actual reference on other
cells.

\section{C3: an exact finite convex extension, with no intermediate
operator}\label{c3-an-exact-finite-convex-extension-with-no-intermediate-operator}

Choose a rational number \(M\) at least the global Lipschitz constant of
\(f\) on \(\sigma\). Define on all of \(Q\)
\[ \widehat f(x)=\inf_{y\in\sigma}\{f(y)+M\|x-y\|\}. \] It is finite: a
fixed \(y_0\in\sigma\) gives an upper bound, and Lipschitzness and the
triangle inequality give the lower bound \(f(y_0)-M\|x-y_0\|\). For
\(x\in\sigma\), every summand is at least \(f(x)\) and the choice
\(y=x\) gives equality. The triangle inequality inside the infimum
proves global \(M\)-Lipschitzness. Convexity follows by taking
approximate minimizers at two points, using their convex combination in
\(\sigma\), the convexity of \(f\), and the triangle inequality for the
norm, and then letting the approximation tolerance go to zero.
Attainment of the infimum is unnecessary and is not asserted.

For a nearest point \(y\in\sigma\), \[
\widehat f(x)\leq f(y)+M\|x-y\|
\leq\ell(x)+C+(M+\|a\|)\operatorname{dist}(x,\sigma)
\leq\Gamma(x)+C
\] provided the integer \(\kappa\) is also at least \(H(M+\|a\|)\).
Increasing \(\kappa\) preserves C2. Set
\[ F(x)=\max\{\widehat f(x),\Gamma(x)-C\}. \] Then \(F\) is finite
continuous convex on all of \(Q\), \(\Gamma-C\leq F\leq\Gamma+C\), and
\(F=f\) on the entire \(\sigma\). The last equality uses
\(\Gamma|_\sigma=\ell\) and \(f\geq\ell-C\). This maximum is a maximum
of continuous convex functions, with no piecewise-affine claim about
\(\widehat f\) or \(F\). No Monge--Ampere measure or mixed operator is
applied to this intermediate function. This completes C3.

\section{C4: full-affine primal majorants, including real heights and
boundary
fibers}\label{c4-full-affine-primal-majorants-including-real-heights-and-boundary-fibers}

We reconstruct the full-affine mechanism rather than citing a grant
label. Let \(Q\) have positive dimension; the point case is immediate. A
finite convex real PA function \(\Gamma\) on \(Q\) is the maximum of its
finitely many actual affine pieces. Each full-dimensional piece is a
global affine support by convex secants; lower-dimensional pieces are
their continuous traces. Thus \[ \Gamma(x)=\max_j(a_j\cdot x+b_j),\qquad
h(v)=\Gamma^\infty(v)=\max_j a_j\cdot v. \] Let
\(D=\operatorname{conv}\{a_j\}\). This compact polytope need not be
rational. The coercive inequality \(h(v)>0\) for nonzero \(v\) implies
that zero is in its interior, by nearest-point separation if it were
not. Both \(\Gamma-h\) and \(F-h\) are bounded on all of \(Q\).

Define the actual conjugate \(F^*(p)=\sup_x(p\cdot x-F(x))\). Its finite
domain is exactly \(D\). For \(p\in D\), \(p\cdot x\leq h(x)\) bounds
the supremum uniformly above, while evaluation at a fixed \(x\) bounds
it below. For \(p\notin D\), separation gives \(v\) with
\(p\cdot v>h(v)\), and \(x=t v\) makes the supremum infinite. On \(D\)
the conjugate is convex, lower semicontinuous, finite, and bounded.

The needed continuity includes boundary fibers. More generally a finite
lower semicontinuous convex function \(g\) on a compact polytope is
continuous relative to that polytope. It is bounded above by the largest
value at its finitely many vertices. At \(p\), choose \(r>0\) below
every inactive facet slack divided by the corresponding row norm. If
\(q\) is near \(p\) in the polytope, the point \(z=p+r(q-p)/\|q-p\|\)
stays in it: active rows do not increase, and inactive rows have the
chosen margin. Convexity gives \[ g(q)\leq g(p)+(\|q-p\|/r)(B-g(p)). \]
Together with lower semicontinuity this proves continuity at every
boundary point as well as in the interior. Work in the relative affine
hull for lower-dimensional polytopes; the single-point case is
immediate. This proof does not assume rational facets of \(D\).

By continuity of \(F\), rational spatial samples compute the same
supremum defining \(F^*\). For each \(p\in D\), choose a rational \(s\)
whose value \(p\cdot s-F(s)\) is within \(\eta/2\) of \(F^*(p)\).
Continuity of the latter on \(D\) and of the displayed affine function
makes the error less than \(\eta\) on a neighborhood of \(p\).
Compactness supplies a finite set of rational samples \(A\) such that
\[ F^*(p)-\eta\leq g_A(p):=\max_{s\in A}(p\cdot s-F(s))\leq F^*(p)
\quad(p\in D). \] Append any prescribed finite rational contact set.
Enlarging \(A\) preserves both inequalities. The sample heights remain
the \emph{actual real numbers} \(F(s)\); they are not rounded.

The recession epigraph has finitely many exact rays with rational
spatial projections. Indeed the rational recession cones of the finite
carrier of \(\Gamma\) cover \(Q\): along a ray a subsequence lies in one
carrier cell, and its defining inequalities show that the direction lies
in that cell's recession cone. On each such cone \(h\) is linear. Refine
them into a finite rational pointed fan, by a rational coordinate
orthant subdivision if necessary, and use their rational generators
\(r_k\). Then
\[ \operatorname{epi}h=\operatorname{cone}\{(r_k,h(r_k)),(0,1)\}. \] For
a direction in one cone, a nonnegative generator representation gives
exactly its height because \(h\) is linear there; the vertical ray
accounts for any additional height. The spatial generators are rational,
but their heights may be irrational. This is the exact recession cone,
not a rationalized slope surrogate.

Form the primal polyhedron
\[ E=\operatorname{conv}\{(s,F(s)):s\in A\}+\operatorname{epi}h. \] It
is closed: the first set is compact, the cone is closed, and a
convergent sum sequence has a convergent compact summand and hence a
convergent cone summand. It is a finite real polyhedron by the
elementary elimination argument. Its spatial projection is all of \(Q\),
since the projection of \(\operatorname{epi}h\) is all of \(Q\); every
fiber contains an upward ray.

Convex secants and bounded \(F-h\) show \(F(x+v)\leq F(x)+h(v)\):
compare the unit secant along \(v\) with arbitrarily long secants and
pass to infinity. Convex combinations of the sample lifts followed by
adding a point of \(\operatorname{epi}h\) therefore remain in
\(\operatorname{epi}F\). Thus \(E\subseteq\operatorname{epi}F\). Every
fiber has a finite attained bottom, defining a finite convex function
\(P\) with \(E=\operatorname{epi}P\). It is continuous by the local
convex bound proof, since its domain is the full affine space.

Direct optimization over the hull coefficients and cone gives
\[ P^*(p)=g_A(p)\quad(p\in D),\qquad P^*(p)=+\infty\quad(p\notin D). \]
For \(p\in D\) the supremum over the cone is zero; outside \(D\) it is
infinite along a separating recession ray. The double-conjugate identity
proved from full-space supports in the foundations now gives
\[ F\leq P\leq F+\eta\quad\hbox{on all of }Q. \] At every retained
sample, the lift belongs to \(E\), so \(P(s)\leq F(s)\); the lower
inequality gives equality. This proves finite-contact exactness and a
global uniform error, not merely convergence on a bounded set.

We must check rational walls despite irrational heights. Every lower
graph face of \(E\) is exposed by a functional with negative vertical
coefficient. Full spatial projection rules out vertical facets, and any
graph face is the intersection of lower facets, so a positive sum of
their exposing normals has negative vertical coefficient. Its exposed
hull summands are a subset of the sample points, and its exposed cone
summands are a subset of the exact rays with zero exposing value. The
vertical ray is not among them. On the graph face, projection is an
affine isomorphism, since that face lies on one affine graph. Its
projected cell is consequently
\[ \operatorname{conv}\{\hbox{rational sample locations}\}+\operatorname{cone}\{\hbox{rational spatial rays}\}, \]
a rational polyhedron by the elementary projection proof. Intersections
project to common faces: the graph over a common projected point is
unique, and the affine graph projection on each face preserves its face
lattice. The finite projected lower-face complex covers \(Q\). On it
\(P\) is real affine and continuous.

Finally this is the genuine scalar-extension notion of real PA, rather
than a collection of arbitrary irrational-wall formulas. The coefficient
tuples for continuous affine functions on this rational carrier satisfy
rational matching equations, obtained by testing equality on rational
affine bases of shared faces. Their solution space has a rational basis
by Gaussian elimination. Clear denominators in each basis tuple to make
its slopes integral on the fixed lattice and its constants rational.
Every actual real coefficient tuple is a real linear combination of
these rational PA basis functions. Thus \(P\) belongs to the required
scalar real PA space. The same argument applies to \(\Gamma\). No
generic closure of real PA under an arbitrary maximum with irrational
switching walls has been assumed.

The recession of \(E\) is exactly \(\operatorname{epi}h\). Inclusion of
this cone is immediate. For the reverse inclusion, divide a recession
sequence by its parameter; the bounded hull part disappears and
closedness leaves its direction in the cone. Equivalently use the
homogenized finite representation. Hence \(P^\infty=h\). Ordinary
full-space convexity also gives local nonnegative divisor inequalities
on any positive balanced test cycle: subtract a supporting affine
function which agrees with the cell's tangent gradient, then each
incident inward residual slope is nonnegative and the balanced affine
contribution vanishes. This justifies the PC property when it is
required on such a test cycle; it does not manufacture a mixed measure
on an unbalanced cell.

An exact one-dimensional illustration retains an irrational height. Take
\(F(x)=\sqrt{1+x^2}\) and reference \(\Gamma(x)=|x|\). The second
derivative of \(F\) is \((1+x^2)^{-3/2}>0\), and \(0<F-\Gamma\leq1\).
The two samples at \(x=\pm1\) have actual height \(\sqrt2\). Their
primal hull plus \(\operatorname{epi}|v|\) gives
\[ P(x)=\sqrt2+\max(0,|x|-1). \] Its spatial walls are exactly the
rational points \(-1,1\), its sample contacts are exact, and its
recession is \(|v|\). For \(r=|x|\leq1\), \(P^2-F^2=1-r^2\geq0\); for
\(r\geq1\), \(P^2-F^2=2(\sqrt2-1)(r-1)\geq0\). Both functions are
positive. Also \(P-(\sqrt2-1)\) equals \(1\) in the first region and
\(r\) in the second, each at most \(F\). Thus \(0\leq P-F\leq\sqrt2-1\)
globally, with equality at zero. The exact finite program checks these
comparisons in \(\mathbb Q(\sqrt2)\) on its specified rational inputs;
the displayed inequalities prove the whole-line assertion. No rational
height approximation is involved.

Restrict the rational carrier of \(P\) to \(\sigma\) by intersection.
The resulting rational carrier is finite, the restriction is convex real
PA, and \[ f\leq P|_\sigma\leq f+\eta. \] Prescribed rational contacts
are exact because \(F=f\) on \(\sigma\). On
\(\operatorname{rec}(\sigma)\), \(r(v)=0\), so \(h(v)=a\cdot v\), the
actual recession of \(\ell\). Bounded \(f-\ell\) has that same recession
along every feasible ray. This verifies the actual same cell recession
without altering the reference. C4 is complete.

\section{C5: decreasing cell approximants and the limits of
contacts}\label{c5-decreasing-cell-approximants-and-the-limits-of-contacts}

Choose \(\eta_n=2^{-n}\) and unshifted majorants \(p_n=P_n|_\sigma\)
with \(0\leq p_n-f\leq\eta_n\). They may all preserve the same
prescribed finite rational contact set. Define \[ g_n=p_n+4\eta_n. \]
Then \(4\eta_n\leq g_n-f\leq5\eta_n\), while
\[ g_n-g_{n+1}\geq4\eta_n-5\eta_{n+1}=\tfrac32\eta_n>0. \] Thus \(g_n\)
decreases uniformly to \(f\) on the whole cell, remains real PA on a
rational carrier, and has the same recession. The shifted functions do
not retain exact contacts; the unshifted ones do. On a point cell the
construction reduces to constants and all assertions are immediate. This
completes all five units of working296.

Finite contacts do not imply equality of two separately constructed
traces on an entire shared positive-dimensional face. For a concrete
finite illustration, the function \(t^2\) on \([0,1]\) has rational-mesh
chord majorants. Two sufficiently fine rational meshes containing any
prescribed finite rational contact set give exact agreement with \(t^2\)
at those contacts, uniform errors as small as requested, but different
traces between mesh points if the meshes differ. The maximum interval
error of a chord is its interval length squared divided by four. This
example proves the logical insufficiency of finite contact constraints
for entire-face trace matching. Nor do separate convex supports on
different incident cells ensure a single ambient vector satisfying all
their support inequalities. The simultaneous direction-cost argument of
L3 applies to the original printed-PC function; it has not been proved
for arbitrarily chosen independent PA approximants.

\section{B1--B3: the unbalanced orthant boundary
diagnostic}\label{b1b3-the-unbalanced-orthant-boundary-diagnostic}

On the positive open quadrant, the Hessian quadratic form of
\(\sqrt{xy}\) is \[ -\frac{(y\,dx-x\,dy)^2}{4(xy)^{3/2}}. \] Thus the
geometric mean is concave there. Continuity and inward approximation
extend its concavity to the closed orthant. Alternatively, the concavity
inequality follows directly by squaring and the two-term Cauchy--Schwarz
inequality. Hence \(1-\sqrt{xy}\) is convex and continuous on the closed
orthant, and so is its maximum with zero. Adding \(x+y\) proves
continuity and convexity of \(f\). Its difference from \(\gamma\) lies
in \([0,1]\). The slope of \(\gamma\) is strictly positive on every
nonzero orthant recession direction. Dividing the bounded difference by
the ray parameter shows \(f^\infty(v)=v_1+v_2\). These are the full B1
assertions.

Every continuous finite real PA function on a convex cell is globally
Lipschitz with the maximum of the norms of its finitely many intrinsic
gradients. For any segment, each carrier cell cuts it in a closed
interval, possibly empty or a point. The finitely many interval
endpoints partition the segment into finitely many subsegments on which
one affine formula applies. Sum the affine Lipschitz bounds over these
subsegments, using continuity at their endpoints. This argument allows
actual real coefficients and even real walls; rational walls are not
needed for the diagnostic. Apply it also to \(a-\gamma\), obtaining some
finite constant \(L\).

Let \(z_n=(n^2,0)\) and \(w_n=(n^2,n^{-2})\). Their Euclidean distance
is \(n^{-2}\), while \((f-\gamma)(z_n)=1\) and \((f-\gamma)(w_n)=0\). If
\(e=\|a-f\|_\infty\) is finite, the triangle inequality gives
\[ 1\leq2e+L n^{-2}. \] Letting \(n\to\infty\) proves \(e\geq1/2\). If
\(e\) is infinite the assertion is immediate. The real PA function
\(a=\gamma+1/2\) attains exactly \(1/2\), so the floor is sharp. This
verifies all of B2 on the continuum; bounded exact computations in the
accompanying program illustrate the limiting inequality but do not
replace it.

The same pair has \(f(z_n)-f(w_n)=1-n^{-2}\) at distance \(n^{-2}\), so
\(f\) is not uniformly continuous and is not globally Lipschitz on the
orthant. At a boundary point \((x,0)\) with \(x>0\), the inward \(y\)
directional quotient equals \(1-\sqrt{x}/\sqrt t\) for small \(t\) and
tends to \(-\infty\). A finite convex extension to an open neighborhood
would be locally Lipschitz by the foundational cube argument and cannot
have this derivative. This explains why the Lipschitz input to C3 cannot
be deleted.

A single positive weighted orthant is unbalanced along either axis:
there is only one inward primitive normal and no opposite contribution.
Its boundary star lies in a proper halfspace. The positive-flux
contradiction of L2 therefore does not apply. This verifies B3 and
precisely locates the diagnostic outside the main theorem. It supplies
no original balanced counterexample and has no mixed-operator
interpretation in this paper.

\section{Scope, source qualifications, and the remaining original
task}\label{scope-source-qualifications-and-the-remaining-original-task}

The original task asks for classifications over balanced complexes and
admissible reference/obstacle domains, with additional contact and
operator conclusions. The present results verify the proposed
finite-geometry Lipschitz theorem, its single-cell consequence, and its
boundary diagnostic. They do not establish membership in a global joint
rational-carrier regularization domain on an arbitrary complex with
positive-dimensional overlaps.

For an admitted actual positive-growth real PA PC reference \(\gamma\),
R230 supplies continuity of the actual envelope over \emph{all} bounded
continuous PC competitors. Its upper-only competitor envelope equals its
two-sided bounded one: for an upper-only competitor \(q\leq\gamma+u\),
replace it by \(\max(q,\gamma-\|u\|_\infty)\). The maximum preserves
continuous PC by the PCC inequality, stays below the obstacle, is
bounded on both sides relative to \(\gamma\), and dominates \(q\). This
operation proves the envelope equality in the PC domain. It does not
assert rational-carrier real PA closure for arbitrary maxima.

Such an actual envelope therefore satisfies the hypotheses of Theorem 1
and, cellwise, Theorem 2. Before applying any completed mixed operator,
however, one must place it in a genuine global domain of compatible
uniform PA approximation. Separate cell approximations, finite contacts,
and original-function supports do not supply that missing step. Entire
shared-face equality and common ambient supports for the PA approximants
remain unresolved. We claim no compatible-extension lemma beyond the
exact single-cell theorem proved here.

The qualified grants retain their original restrictions. A224 is a
finite actual-PA strict-growth obstruction, not arbitrary real-pair
classification. B225's finite globally embedded positive-dimensional
noncompactness is a literal finite-global statement, not a prohibition
on all abstract compact atlases. S291's compact maximum result belongs
to its whole boundaryless balanced-atlas domain; only its local
foundations, independently reconstructed here, are used in the main
proof. S262's graph regularization and contact sufficiency do not give
higher-dimensional all-face gluing. S272's real signed and zero
intermediate weights and its zero-mass behavior are retained; finite PA
contact necessity is not non-PA contact sufficiency. Uniform
global-vague control, continuity for ordinary weak nets, and uniform
weak or bounded-Lipschitz control on a fixed relative bound \(C\) are
distinct statements. Global weak/BL uniformity without a fixed \(C\) is
not inferred.

S291 is accepted as a whole mathematical grant with both correction and
reader records. The historical native218 original-source FAIL, combined
FAIL, and retrospective-history INCONCLUSIVE remain. Correcting the
literal LF endpoint of the whole BBGS commutation proof to 670 does not
retroactively turn those reports into PASS. Original-source rational
coefficients and weights are not relabeled as real; the real extension
is supplied separately. Actual nativeFINAL and later saved expansions
remain separate historical objects. We have read only the bound records,
not unbound historical reports or runners, and do not certify those past
runtime events.

No S282, native222/223 outcomes, papers295/298, working293294,
privateACTIVE, or unbound feedback is used. No source or supplied runner
was modified or executed. The four frozen original files are the only
primary sources. Hash and whole-unit LF metadata appear in the ledger.
All statements here are drafts for mathematical checking: successful
finite tests and a local TeX build establish reproducibility, not
independent mathematical acceptance. The original P05.1 and P05.2
classifications remain OPEN; candidate=false.

\section{Reproducibility and proof
coverage}\label{reproducibility-and-proof-coverage}

The companion code performs bounded exact-arithmetic sanity checks for a
spanning balanced star, a nearest-distance example, a uniform inset on
an unbounded face, the saturated covolume ratio, dyadic shifts, and the
orthant obstruction inequalities. These are exhaustive finite results
only over the explicitly enumerated inputs. The general statements above
are mathematical proofs and do not rely on those finite enumerations. No
resource-heavy computation was launched; no supervised compute facility
was supplied or needed for these proofs.

The build uses local Pandoc and pdfLaTeX with shell escape disabled and
captures the first build, its warnings, and subsequent unchanged-source
passes. The programs, bounded inputs, outputs, environment report,
source-reading log, artifact hashes, full notes, and rerun commands are
retained in the fresh artifact directory. No provider nesting, network
request, or agent delegation occurs. The visible execution environment
is recorded without exposing authentication or runtime configuration.
The operator's requested model identity is task metadata, not a
separately verifiable backend certificate.

The coverage is: L0 real perspective density and inherited-cell
convexity; L1 distance and all-slack insets; L2 real covolume flux and
two-sided generator coefficients; L3 finite costs, sublinear envelope,
and entire-cell supports; L4 all-face induction and linear constants; L5
scoped regularization input. C1 distance; C2 coercive auxiliary
reference; C3 exact full-affine extension; C4 actual-dual full-affine
primal majorants; C5 decreasing shifts and finite-contact limits. B1
convexity and recession; B2 the sharp continuum error floor; B3 the
boundary-domain diagnosis. None of these coverage labels includes
arbitrary shared-face gluing or either complete original classification.

\section{Quote-free primary
provenance}\label{quote-free-primary-provenance}

BMPS is the frozen arXiv v1 at \url{https://arxiv.org/pdf/2603.08664v1},
submitted 9 March 2026 with manuscript date 10 March 2026. Its PDF
SHA-256 is:

\texttt{5cf7e68ee9113f1dc813f9aefa5f9ca72cba02fc4fa688f1a61e4b7f0d1223e7}.

Its frozen text SHA-256 is:

\texttt{c4475ad5cf4e8fce3eb6add155fdb81e352cbe079ca0d7646553f8a3e8ad9271}.

Immediate whole units used for definitions and scope are 2.1--2.4,
2.18--2.31, 2.33--2.34, 3.1--3.9, 3.14, 3.22--3.23, 4.15 and 6.1--6.9.
Their precise literal-LF read ranges and unused subarguments are
recorded in the ledger. Neither the original rational coefficient
category nor a printed automatic-continuity assertion is imported as an
actual-real or raw-PC theorem.

BBGS is the frozen publisher version at
\url{https://link.springer.com/content/pdf/10.1007/s00209-021-02869-6.pdf},
Math. Z. 301 (2022), 1631--1674, published 24 January 2022. Its PDF
SHA-256 is:

\texttt{c2817a01e51ac3781e94005cb587c185ac7e6cc5fe654feda6c5341eaf6a5846}.

Its frozen text SHA-256 is:

\texttt{87f8f8d2d51e523c4ecf5f8429875da16994abbe74132ce05886d1782e90b635}.

Immediate whole units read are 2.1--2.3, 2.14--2.15, 3.1--3.9, 3.12,
3.15 and 3.18--3.19. Whole 3.15 occupies literal LF604--670, including
the ending equation and proof. Its actual-real weights, Euclidean unit
normals and negative product sign are kept separate from BMPS's
primitive-lattice-normal convention. The covolume identity in L2
provides the explicit conversion.

The primary PDFs are hash-verified; the mathematical unit reads use the
frozen extracted texts. No claim of reading every PDF page or every line
of either primary is made. All immediate units essential to the fresh
applications were read, and recursive external results mentioned by a
source were either unused or replaced by the elementary proofs above.
The six qualified supplied grants and all their bound annotations were
read whole. Long excerpts and pixels remain local only. No third-party
full text is embedded in the paper or its public ledger.

\end{document}
